\section{The recurrent simplex and the cyclic specialization}
\label{sec:r73-normal-form-supplement}

\begin{lemma}[A lazy finite-chain projection]\label{lem:r73-lazy-convergence}
Let $T$ be a finite row-stochastic matrix with strictly positive diagonal.
Let $C_1,\ldots,C_k$ be its closed communicating classes. There is a
stochastic idempotent $E=\lim_{n\to\infty}T^n$ whose image on the probability
simplex is the simplex spanned by the stationary rows $\mu_i$ supported on
the respective $C_i$. If a row-stochastic matrix $Q$ preserves every $C_i$,
then $EQE=E$. If $T$ is a convex combination of matrices from a multiplicative
semigroup $\mathcal S$ that contains the identity, then $E$ belongs to the
closed convex hull of $\mathcal S$.
\end{lemma}
\begin{proof}
Consider first one closed class $C$. Irreducibility and the positive
diagonal imply that some power $K$ of the restricted matrix has every
entry positive: choose a path between each ordered pair, and pad the
finitely many path lengths to their maximum using diagonal self-loops.
Put $\eta=\min_{i,j}K_{ij}>0$ and $n_C=|C|$. Then
\[
 K=n_C\eta J+(1-n_C\eta)K_0,
\]
where all rows of $J$ are uniform and $K_0$ is row-stochastic, with the
second term omitted if $n_C\eta=1$. For probability rows $p,q$,
\[
 \|pK-qK\|_1\le(1-n_C\eta)\|p-q\|_1.
\]
A stationary row for $T|_C$ exists by taking an accumulation point of the
Ces\`aro averages $n^{-1}\sum_{j=0}^{n-1}pT^j$ in the compact simplex:
their difference from their right translates by $T$ tends to zero.
Call it $\mu_C$. The contraction shows it is unique and that all rows
converge geometrically to $\mu_C$ along the chosen power. The remaining
finitely many residues of the exponent converge to the same row because
$\mu_CT=\mu_C$. This proves convergence on each closed class.

Let $H$ be the complement of their union. Every vertex in $H$ has a
directed path to a closed class. Since the graph is finite, there are
an integer $q$ and a number $\zeta>0$ such that, starting from any vertex
in $H$, the probability of leaving $H$ by time $q$ is at least $\zeta$.
Paths can again be padded using the positive diagonal. In a permutation
of coordinates, write
\[
 T=\begin{pmatrix}A&B\\0&C\end{pmatrix},
\]
where $A$ is indexed by $H$ and $C$ is block diagonal over the closed
classes. Then $\|A^n\|_\infty\le(1-\zeta)^{\lfloor n/q\rfloor}$, while
$C^n\to E_C$, the block projection with all rows in $C_i$ equal to
$\mu_i$. The upper-right block of $T^n$ is
\[
 \sum_{j=0}^{n-1}A^jBC^{n-1-j}.
\]
This converges to $\sum_{j=0}^\infty A^jBE_C$: for a fixed cutoff the
finitely many factors $C^{n-1-j}$ converge, and the tail is uniformly
bounded by a convergent geometric series. Thus
\[
 E=\begin{pmatrix}0&\sum_{j=0}^\infty A^jBE_C\\0&E_C\end{pmatrix}
\]
exists. As a limit of stochastic matrices it is stochastic. Every row
is a convex combination of the $\mu_i$, each of which occurs as a row
of $E$ on its own class. Their disjoint supports give affine independence,
and $\mu_iE=\mu_i$ proves $E^2=E$.

If $Q$ preserves $C_i$, the row $\mu_iQ$ is a probability law on $C_i$,
so $(\mu_iQ)E=\mu_i$. Extend over the convex combinations constituting
the rows of $E$ to obtain $EQE=E$. This does not require $\mu_iQ=\mu_i$.
Finally, every $T^n$ is a convex combination of products from
$\mathcal S$, which remain in $\mathcal S$. Taking the limit gives
the last assertion.
\end{proof}

\subsection{Cyclic modules and the retained positive-amplitude model}
The next statement identifies exactly where the representation-theoretic
bound is sufficient for the periodic family; it is not asserted for every
response on a cyclic group.

Let $G=C_Q$ and let the distinct nontrivial character pairs occurring in a
real $G$-module $V$ have orders $q_1,\ldots,q_r$ dividing $Q$; equal orders
are allowed for different pairs. Assume each complex-conjugate pair occurs
once, the trivial character occurs once, and no other characters occur.
For this assertion a real character of order two, if present, is counted
once. For a partition $\mathcal P$
of $\{1,\ldots,r\}$, put
\[
 \Psi(\mathcal P)=\sum_{B\in\mathcal P}\operatorname{lcm}(q_j:j\in B).
\]

\begin{proposition}[Cyclic permutation-module minimum]
\label{prop:r73-cyclic-module}
If $r\ge1$, the least size of a permutation module containing $V$ is
\begin{equation}\label{eq:r73-cyclic-module}
 \nu_{C_Q}(V)=\min_{\mathcal P}\Psi(\mathcal P).
\end{equation}
For the nondegenerate positive-amplitude periodic response of
Section~\ref{sec:r72-periodic}, the cycle-mixture construction there
attains this bound with a probability decoder. Consequently
$\kappa_G(b)=\nu_G(V_b)$ for that family, and
\eqref{eq:r73-memory-formula} gives its retained partition minimum plus
the pure-start alternative. A constant response has $\nu=\kappa=m=1$.
\end{proposition}
\begin{proof}
A transitive $C_Q$-set is a cycle of length $\ell$ dividing $Q$.
Its real permutation module contains exactly the characters whose orders
divide $\ell$, once each after the conjugate pairs are combined over
$\mathbb R$. This can be verified directly with the vectors
$j\mapsto\exp(2\pi imj/\ell)$, $0\le m<\ell$, which are a basis and
diagonalize cyclic translation. Given any collection of cycles whose
module contains $V$, assign each required nontrivial character to one
cycle that contains it, and discard cycles to which no character was
assigned. For the nonempty index block $B$ assigned to a cycle, its
length is a multiple of every $q_j$, $j\in B$, and therefore at least
their least common multiple. The total number of points is at least
$\Psi(\mathcal P)$ for the induced partition. Conversely, cycles of
the lengths in any $\Psi(\mathcal P)$ contain every required character
and at least one trivial character. This proves \eqref{eq:r73-cyclic-module}.

For the indicated positive-amplitude response, the construction in
Section~\ref{sec:r72-periodic} uses one such cycle for each block,
chooses its initial block weight $W_B/\Gamma$, where
$W_B=\sum_{j\in B}w_j$ and $\Gamma=\sum_jw_j\le1$, and uses the binary
decoder
\[
 d_{B,n}=\frac12\left[1+\frac{\Gamma}{W_B}
       \sum_{j\in B}w_j\cos(2\pi n a_j/q_j+\varphi_j)\right].
\]
Each cosine lies in $[-1,1]$ and the coefficients have total $\Gamma\le1$,
so every decoder lies in $[0,1]$. It realizes the response with exactly
the cycle count in \eqref{eq:r73-cyclic-module}. The opposite inequality
is \eqref{eq:r73-module-bounds}. The separate initial label is charged
exactly as in Theorem~\ref{thm:r73-memory-formula}; when the residual
orbit itself is cheaper, its pure-start machine is used.
\end{proof}

The cyclic conclusion uses the amplitudes, frequencies, and normalization
actually specified in the retained theorem. Opposite frequencies must be
combined, zero amplitudes removed, and order-two modes treated as real
one-dimensional modes before applying the module statement. It does not
replace the original hypotheses by a mere list of eigenvalue orders.
